\documentclass[preprint,12pt]{elsarticle}

\usepackage[T1]{fontenc}
\usepackage[utf8]{inputenc}
\usepackage{amsmath,amssymb}
\usepackage{booktabs}
\usepackage{graphicx}
\usepackage{multirow}
\usepackage{url}
\usepackage{hyperref}
\usepackage{enumitem}
\usepackage{float}
\usepackage{natbib}

\journal{Computers and Education: Artificial Intelligence}

\begin{document}

\begin{frontmatter}

\title{Analyzing Generalizability and Optimizing Budget Forcing for Test-time Scaling across Open-source LLMs}

\author[inst1]{Van Hong Hanh}
\author[inst1]{Phan Hong Hanh}
\address[inst1]{University of Technology, Ho Chi Minh City, Vietnam}

\begin{abstract}
This report studies budget forcing from \citet{muennighoff2025s1} as a simple test-time scaling mechanism for reasoning-oriented language models. The project reproduces the core decoding intervention implemented in this repository, evaluates it on mathematical reasoning benchmarks, and analyzes whether increasing the allowed thinking budget improves downstream accuracy under constrained resources. In parallel with the reproduction effort, the report situates budget forcing within the broader landscape of inference-time enhancement methods, including retrieval augmentation and reasoning-oriented post-training. The current artifact is structured to support a narrow and reproducible reproduction: one primary benchmark, one main model configuration, a small sweep over \texttt{n\_wait}, and a trigger-phrase ablation. Final result paragraphs, tables, and figures are designed to be filled directly from the JSON outputs produced by the experiment pipeline.
\end{abstract}

\begin{keyword}
test-time scaling \sep budget forcing \sep reasoning language models \sep reproduction study \sep mathematical reasoning
\end{keyword}

\end{frontmatter}

\section{Introduction}

Recent reasoning-focused language models have renewed interest in test-time scaling, where extra inference-time computation is used to improve answer quality rather than expanding training data or model parameters. The s1 project argues that strong reasoning gains do not necessarily require a complex training stack, large proprietary datasets, or opaque systems. Instead, it combines a compact supervised fine-tuning dataset, called s1K, with a decoding-time intervention named budget forcing \citep{muennighoff2025s1}.

This repository focuses on the budget forcing component because it is narrow, inspectable, and experimentally accessible. In contrast to larger reinforcement-learning-based systems such as DeepSeek-R1 \citep{deepseek2025r1} or large-scale inference-oriented models such as Kimi k1.5 \citep{kimi2025k15}, budget forcing modifies generation at inference time by either extending or truncating the model's reasoning trace. That makes it a suitable target for a graduate reproduction project with limited compute.

The present work has two goals. First, it audits and executes the existing implementation in this repository to establish whether the reported budget-forcing behavior can be reproduced under a modest experimental setup. Second, it analyzes the mechanism and its limits through the lens of reproducibility: how much of the observed behavior depends on model choice, trigger design, benchmark selection, and hardware availability.

The main contributions of this report are as follows. We provide a structured summary of the s1 method and its evaluation metrics, instantiate an end-to-end experiment pipeline for budget forcing on math reasoning tasks, and document a paper-style reproduction workflow that separates stable narrative content from result-dependent content. This separation is useful because the repository was already close to feature-complete, while the missing work was concentrated in execution, reporting, and integration.

The remainder of the paper is organized as follows. Section~\ref{sec:method} summarizes the s1 method and the local implementation. Section~\ref{sec:metrics} explains the three evaluation metrics used to analyze test-time scaling. Section~\ref{sec:analysis} discusses strengths, limitations, and research gaps. Section~\ref{sec:related} connects budget forcing to adjacent work on reasoning models and inference-time augmentation. Section~\ref{sec:experiments} documents the experimental setup and provides a template for presenting the observed results. Section~\ref{sec:conclusion} concludes with practical lessons and future directions.

Test-time scaling has become a central topic in recent reasoning-focused large language model research because it offers a path to stronger performance without changing the base pretraining corpus at inference time. Instead of relying only on larger models or more supervised data, test-time scaling methods spend additional compute during generation, typically through longer reasoning traces, sampling, verification, or structured control policies. Recent public discussion around reasoning models has intensified because systems such as OpenAI's o1 demonstrated strong reasoning behavior without fully disclosing the recipe behind that behavior.

The s1 framework is important in this context because it argues that competitive reasoning behavior can emerge from a comparatively simple recipe: a small, carefully curated dataset of reasoning traces and a decoding-time intervention called budget forcing \citep{muennighoff2025s1}. Rather than treating reasoning depth as a fixed model property, budget forcing manipulates the generation process directly. When the model tries to stop too early, the decoder can suppress termination and append a trigger phrase such as ``Wait'', effectively encouraging more internal computation. Conversely, when a hard budget is reached, the decoder can force a transition to the final answer segment.

This project examines that claim from a reproduction perspective. Our repository focuses on the decoding-time mechanism rather than a full-scale retraining pipeline. The implemented code already contains a budget forcing decoder, an evaluation entrypoint, basic scripts for baseline and budget forcing sweeps, and notebooks for reproduction and trigger ablation. The main research task now is to turn these components into validated experimental artifacts and a coherent analysis.

The report makes three practical contributions. First, it explains the s1 method and the role of budget forcing in accessible terms. Second, it documents the current repository implementation and its evaluation design. Third, it frames reproduction as an engineering problem with explicit environmental constraints, especially for local execution on non-CUDA hardware. This is important because negative or partial reproduction results can still be scientifically useful when they clarify what parts of a claimed method are robust and what parts depend on hidden infrastructure assumptions.

The remainder of this report is organized as follows. Section~\ref{sec:method} summarizes the s1 pipeline and the implemented budget forcing mechanism. Section~\ref{sec:metrics} explains the control, scaling, and performance metrics used to judge test-time scaling behavior. Section~\ref{sec:analysis} discusses the strengths, limits, and research gaps that motivate our reproduction. Section~\ref{sec:related} positions the work relative to adjacent reasoning and inference-time augmentation methods. Section~\ref{sec:experiments} describes the repository experiments and records where result integration depends on successful execution. Section~\ref{sec:conclusion} concludes with key observations and next steps.
\section{Method Summary}\label{sec:method}

\subsection{s1 and the s1K data curation strategy}

The s1 paper proposes a deliberately simple recipe for eliciting reasoning improvements from an open model \citep{muennighoff2025s1}. Instead of relying on very large curated reasoning corpora, the authors construct a 1,000-example training set, s1K, guided by three properties: quality, difficulty, and diversity. The key claim is that a small but carefully filtered dataset can be sufficient when paired with a model family that already has a usable latent reasoning prior.

Although this repository does not reproduce the full supervised fine-tuning stage, the project plan and codebase inherit the paper's framing. The existing scripts assume that a reasoning-capable checkpoint, such as a DeepSeek-R1 distilled model or a Qwen-based instruction model, can serve as the experimental base for budget forcing without re-implementing the entire training pipeline.

\subsection{Budget forcing as a decoding-time intervention}

Budget forcing modifies generation after the model has been prompted to think step by step. In the implementation under \texttt{experiments/budget\_forcing/decoding.py}, the decoder wraps the model's autoregressive generation loop and intervenes in two ways.

First, enforce-minimum behavior suppresses an early end-of-thinking marker and appends a trigger phrase, usually \texttt{Wait}, when the model attempts to stop before the configured thinking budget has been exhausted. This is intended to induce self-reflection or reconsideration. Second, enforce-maximum behavior injects a final-answer prefix when the model exceeds a designated maximum thinking budget, forcing it to transition from internal reasoning to answer production.

The mechanism is appealing from a systems perspective because it does not require retraining, external tools, or a second verification model. At the same time, the same simplicity creates identifiable risks: the trigger phrase may cause repetitive loops, reasoning quality may degrade as the trace becomes longer, and the intervention may interact differently with different chat templates or end-of-turn markers.

\subsection{Local pipeline in this repository}

The main experiment entrypoint is \texttt{experiments/evaluation/run\_eval.py}. It formats a math-oriented chat prompt, loads a benchmark from Hugging Face datasets, runs the configured model through the budget-forcing decoder, extracts a predicted answer, and computes accuracy by comparing it with the benchmark answer. Shell wrappers under \texttt{experiments/scripts/} expose a baseline run and a sweep over \texttt{n\_wait} values.

This repository therefore supports a direct reproduction pattern: baseline decoding with \texttt{n\_wait = 0}, followed by budget-forcing runs for \texttt{n\_wait = 1, 2, 4}, optionally repeated with an alternative trigger phrase. The outputs are serialized as JSON files in \texttt{experiments/results/}, which then serve as the data source for notebooks, plots, and the final report.
\section{Evaluation Metrics}\label{sec:metrics}

The s1 paper evaluates test-time scaling with three metrics: control, scaling, and performance \citep{muennighoff2025s1}. The repository mirrors this structure in \texttt{experiments/budget\_forcing/metrics.py}.

\subsection{Control}

Control measures whether the method can reliably reach a target compute budget. In the local implementation, control is computed as the fraction of runs whose observed token count lies within a tolerance band around the target budget. This metric matters because a method that occasionally spends more compute is not sufficient; test-time scaling should be deliberate rather than accidental.

In the current repository state, control is only fully meaningful when target budgets and actual thinking-token counts are both logged in a way that can be matched run by run. For the default \texttt{n\_wait} sweep, the code explicitly notes that control may be unavailable or only partially approximated.

\subsection{Scaling}

Scaling captures whether additional inference-time compute translates into higher accuracy. The local implementation estimates this using the slope of a linear fit over accuracy versus compute level, with accuracy converted to percentage points. Positive scaling indicates that larger compute budgets improve performance; negative scaling suggests that longer reasoning traces are harming the model.

This metric is especially important for a reproduction study because it separates two distinct outcomes. A model may achieve a reasonable peak accuracy while still failing to scale with extra compute, or it may scale positively but only over a narrow range of budgets. In both cases, reporting a single maximum accuracy would obscure the real behavior.

\subsection{Performance}

Performance is defined as the maximum accuracy observed across tested compute levels. It is the simplest metric to interpret and the easiest one to compare against prior work. However, it is also the least informative in isolation. A reproduction that reaches a similar peak performance but only at a much higher compute cost, or with unstable trigger behavior, should not be treated as equivalent to the original claim.

Together, these three metrics provide a balanced view of test-time scaling. Control evaluates budget precision, scaling evaluates the direction and efficiency of extra compute, and performance evaluates the best-achieved outcome.
\label{sec:metrics}

The s1 paper emphasizes that test-time scaling should not be judged by accuracy alone. A useful control method must also manipulate compute in a predictable way. The repository mirrors this framing through three metrics implemented in \texttt{experiments/budget\_forcing/metrics.py}.

\subsection{Control}

Control measures whether the realized compute stays close to the intended budget. In the local implementation, this is defined as the fraction of runs whose actual token count falls within a tolerance band around the target budget. The motivation is straightforward: a method that claims to spend more or less test-time compute must actually do so reliably. If the model cannot follow the intended budget, downstream comparisons become hard to interpret.

\subsection{Scaling}

Scaling measures how accuracy changes as more compute is allocated. The repository uses the slope of an ordinary least squares fit over compute levels and accuracies, scaled into percentage units. Positive scaling indicates that additional compute is associated with better performance; negative scaling suggests wasted or counterproductive computation. This metric is especially relevant for budget forcing because the technique is useful only if longer controlled reasoning often helps instead of merely prolonging incorrect traces.

\subsection{Performance}

Performance is the maximum accuracy achieved across the tested compute levels. This metric captures the best outcome obtained from the current compute-control schedule, regardless of whether all intermediate points are monotonic. It is a practical summary for a coursework reproduction because it answers the simplest applied question: what is the strongest accuracy the current setup reached?

\subsection{Interpretation in This Project}

In the present repository state, control may remain partially unavailable until the run configuration explicitly records target budgets and actual reasoning-token counts for each sample. By contrast, scaling and performance can already be computed from accuracy values across different \texttt{n\_wait} settings. This asymmetry is important and should be reported honestly in the experiments section rather than hidden. A partial metric profile is still informative if the missing piece is clearly explained.
\section{Strengths, Limitations, and Research Gaps}\label{sec:analysis}

\subsection{Strengths}

The main strength of budget forcing is its accessibility. Compared with reinforcement-learning-heavy reasoning pipelines, the method is conceptually compact and locally inspectable. This repository reflects that advantage: the core intervention is implemented in a single decoding wrapper, and the evaluation path is understandable without reverse-engineering a large training stack.

Another strength is reproducibility by decomposition. Even when the full s1 training setup is not reproduced, the core claim about inference-time behavior can still be examined through targeted experiments. This makes budget forcing a useful case study for academic courses that want to study reasoning behavior without retraining a frontier-scale model.

\subsection{Limitations}

The simplicity of the method also exposes its limits. First, budget forcing depends on model-specific control tokens and prompt structure. The repository currently assumes reasoning-style chat templates and end-of-turn conventions that may not transfer cleanly across model families. Second, extending the reasoning trace with a fixed trigger phrase may encourage repetitive or low-value continuation rather than meaningful self-correction. Third, the local implementation is benchmarked around math tasks, so conclusions should not be generalized to other domains without further evidence.

From a practical standpoint, infrastructure remains a major constraint. The repository defaults to 4-bit loading with bitsandbytes and is optimized for CUDA-capable hardware. That makes the reproduction pathway more fragile on macOS or CPU-only systems, where a complete run may be infeasible even if the code itself is logically correct.

\subsection{Research gaps}

The project plan identifies trigger optimization as the most feasible extension for the group. This is a sensible target because trigger choice is both experimentally accessible and theoretically underexplored. A short phrase such as \texttt{Wait} may work as a reflection cue, but it is unclear when such a cue promotes verification versus when it merely prolongs the trajectory.

Two additional gaps remain important. The first is anti-repetition control: a system that extends reasoning should ideally also detect unproductive loops. The second is task breadth. If budget forcing is to be treated as a general test-time scaling primitive, it should be evaluated beyond competition-style mathematics, including tasks where factual grounding, tool use, or retrieval matter more than pure symbolic reasoning.
\label{sec:analysis}

\subsection{Strengths of the s1 Framing}

The main strength of s1 is methodological clarity. By reducing the story to a curated reasoning dataset and a lightweight decoding-time controller, the paper makes test-time scaling feel accessible and reproducible compared with opaque proprietary reasoning systems. This is academically valuable because it exposes a concrete mechanism that can be audited, modified, and tested on new models.

Another strength is sample efficiency. The claim that a relatively small reasoning dataset can support strong downstream behavior challenges the assumption that ever-larger supervised corpora are always necessary. Even if exact performance numbers vary across replications, the idea that data quality and inference control can substitute for brute-force scale is an important hypothesis.

\subsection{Limitations Relevant to Reproduction}

The first limitation is infrastructure dependence. Even a ``simple'' reasoning reproduction can become expensive when it relies on large reasoning-capable models, quantization libraries, and benchmark-scale generation. Course-project environments often differ sharply from the CUDA-heavy setups implicitly assumed by many open-source reasoning papers.

The second limitation is domain concentration. The original framing is especially compelling in mathematical reasoning, but that also narrows the evidence base. It remains unclear how much of the observed benefit transfers to non-math settings without additional adaptation.

The third limitation is behavioral instability. Trigger-based reasoning extension can produce useful self-correction, but it can also create repetitive loops or longer incorrect chains. That makes trigger selection and stopping criteria natural ablation targets.

\subsection{Research Gaps for This Project}

This repository is already aligned with one concrete research gap: trigger optimization. The project plan explicitly identifies trigger-phrase ablation as a manageable extension, and the existing script surface already includes an alternative trigger. This is a sensible direction because the trigger is both conceptually central and experimentally cheap to vary.

Another gap concerns evaluation realism under constrained hardware. Many reproduction attempts fail not because the underlying idea is wrong, but because local execution conditions differ from the source setting. Documenting those constraints carefully turns an otherwise frustrating engineering limitation into a useful piece of scientific context.

Finally, there is a broader gap between internal-knowledge scaling and external-knowledge augmentation. The reference RAG survey shows that another major line of inference-time improvement adds retrieved information instead of lengthening reasoning alone \citep{li2025rag}. Comparing these two paradigms sharpens the conceptual contribution of budget forcing: it primarily tries to improve how a model thinks, not what external evidence it sees.
\section{Related Work}\label{sec:related}

\subsection{Reasoning-oriented test-time scaling}

Budget forcing belongs to a broader family of methods that attempt to improve reasoning by increasing or restructuring inference-time computation. The recent wave of reasoning-focused systems includes large proprietary models and open reproductions, but many of them rely on training pipelines or infrastructure that are difficult to reproduce directly. DeepSeek-R1 demonstrates that reinforcement learning can induce self-reflective reasoning behaviors at scale \citep{deepseek2025r1}. Kimi k1.5 similarly explores stronger reasoning through large-scale reinforcement learning and long-context optimization \citep{kimi2025k15}. Against this backdrop, s1 is notable because it argues for a minimal recipe: a compact curated dataset and a decoding intervention rather than a heavy post-training stack \citep{muennighoff2025s1}. Recent systems such as Sky-T1 also pursue open reasoning recipes that reduce dependence on closed-source pipelines \citep{zhang2025skyt1}.

The repository used in this project aligns with that minimal philosophy. It does not attempt to compete with large closed systems directly. Instead, it asks whether the same qualitative trend, namely better answers under controlled increases in test-time compute, can be observed in a small and reproducible setup.

\subsection{Model backbone and evaluation tooling}

The experimental path in this repository relies on open-weight reasoning or instruction-tuned backbones, particularly DeepSeek-R1 distilled checkpoints and Qwen2.5 variants. Qwen2.5 is a relevant backbone family because it provides instruction-tuned models across multiple sizes and explicitly emphasizes long-context and structured-output capabilities \citep{qwenteam2024qwen25}. For evaluation tooling and reproducibility practices, the project also draws on the design philosophy of the Language Model Evaluation Harness, which standardizes benchmark execution and result logging across model backends \citep{gao2024lmharness}.

\subsection{Inference-time augmentation beyond budget forcing}

Inference-time improvement is not limited to longer reasoning traces. Retrieval-augmented generation, for example, supplements the model's internal knowledge with external evidence rather than extending its internal chain of thought. The systematic survey by \citet{li2025ragsurvey} shows how retrieval augmentation can improve reliability, freshness, and domain grounding in educational applications. This contrast is useful for the present report. Budget forcing attempts to extract more value from a model's internal reasoning process, whereas retrieval augmentation changes the information available at inference time. Both approaches increase inference-time effort without retraining the base model, but they solve different failure modes.

For that reason, the comparison between budget forcing and retrieval-augmented generation should be treated as complementary rather than competitive. One extends internal deliberation; the other injects external evidence. A practical future direction is to study whether the two can be combined, for example by retrieving task-relevant context and then allocating a controllable reasoning budget over that context.
\section{Experiments}\label{sec:experiments}

\subsection{Setup}

The repository is configured around a lightweight reproduction protocol. The main benchmark target is \texttt{math500}, the default sample budget is 50 examples for initial runs, and the default \texttt{n\_wait} sweep is \{0, 1, 2, 4\}. The shell wrappers currently default to the \texttt{r1-distill-7B} checkpoint and use 4-bit loading to reduce memory usage.

All reported runs should include the following metadata: model name, quantization setting, hardware, sample count, trigger phrase, wall-clock runtime, and artifact path. These details are essential because the repository is intended as a reproducible coursework artifact rather than a one-off demo.

\subsection{Baseline results}

This subsection will summarize the no-budget-forcing baseline generated from:

\begin{verbatim}
bash experiments/scripts/run_baseline.sh
\end{verbatim}

Insert the resulting JSON artifact path here and summarize the observed baseline accuracy, average thinking-token count, and any notable answer-extraction failures.

\subsection{Budget forcing results}

This subsection will summarize the main \texttt{n\_wait} sweep generated from:

\begin{verbatim}
bash experiments/scripts/run_budget_forcing.sh
\end{verbatim}

Table~\ref{tab:main-results} is a placeholder for the final accuracy-versus-compute summary.

\begin{table}[t]
\centering
\caption{Main budget-forcing results on the primary benchmark. Replace placeholder values with observed numbers from \texttt{experiments/results}.}
\label{tab:main-results}
\begin{tabular}{lccc}
\toprule
Setting & Accuracy & Avg. thinking tokens & Notes \\
\midrule
Baseline (\texttt{n\_wait = 0}) & TBD & TBD & Greedy decoding \\
Budget forcing (\texttt{n\_wait = 1}) & TBD & TBD & Primary trigger \\
Budget forcing (\texttt{n\_wait = 2}) & TBD & TBD & Primary trigger \\
Budget forcing (\texttt{n\_wait = 4}) & TBD & TBD & Primary trigger \\
\bottomrule
\end{tabular}
\end{table}

Figure~\ref{fig:acc-compute} should plot accuracy against compute level once the notebooks have produced the final figure.

\begin{figure}[t]
\centering
\fbox{\parbox[c][5cm][c]{0.85\linewidth}{\centering Placeholder for accuracy-versus-compute curve generated from the experiment outputs.}}
\caption{Accuracy as a function of the budget-forcing compute level.}
\label{fig:acc-compute}
\end{figure}

\subsection{Trigger ablation}

The repository already anticipates an alternative trigger, \texttt{Hmm, let me reconsider}. This subsection should compare the default trigger with at least one alternative and discuss whether the longer continuation appears to promote useful verification or merely repetitive continuation.

\begin{table}[t]
\centering
\caption{Trigger phrase ablation placeholder.}
\label{tab:trigger-ablation}
\begin{tabular}{lccc}
\toprule
Trigger & Best accuracy & Avg. thinking tokens & Qualitative note \\
\midrule
Wait & TBD & TBD & TBD \\
Hmm, let me reconsider & TBD & TBD & TBD \\
\bottomrule
\end{tabular}
\end{table}

\subsection{Error analysis}

The final version of this subsection should include a small number of representative cases in which budget forcing changes the outcome. In particular, the report should highlight at least one case where the extra reasoning budget corrects an initially flawed trajectory and at least one case where it causes a loop or a degraded answer. This is more informative than reporting aggregate accuracy alone because it reveals whether the trigger is inducing genuine self-correction.

\subsection{Pending artifacts from Agent A}

The following items still depend on experiment execution:

\begin{itemize}
  \item baseline JSON artifact path
  \item budget-forcing sweep JSON artifact path
  \item trigger ablation JSON artifact path
  \item final accuracy-versus-compute plot
  \item case-study excerpts for error analysis
\end{itemize}
\label{sec:experiments}

\subsection{Experimental Setup}

The repository exposes a compact experimental workflow. The default configuration targets a reasoning-capable 7B model, the MATH-500 benchmark, and a sweep over \texttt{n\_wait} values of 0, 1, 2, and 4. The evaluation entrypoint is \texttt{experiments/evaluation/run\_eval.py}, while convenience scripts are provided in \texttt{experiments/scripts/}.

At the time of writing, the most important reproducibility variable is hardware. The code supports optional 4-bit loading, but the default path still assumes a runtime environment compatible with the chosen model and quantization stack. Because this repository is being audited from a macOS workspace, the first experimental responsibility is to validate the command path and determine whether full local execution is feasible.

\subsection{Baseline and Budget Forcing Runs}

The planned baseline is a greedy run with \texttt{n\_wait = 0}. The planned budget forcing sweep evaluates \texttt{n\_wait} values $\{0,1,2,4\}$ using the default trigger phrase ``Wait''. A secondary ablation compares that trigger against at least one alternative phrase already wired into the shell script.

Table~\ref{tab:main-results} is reserved for the final numerical results.

\begin{table}[H]
\centering
\caption{Main reproduction results. Fill this table after Agent A produces the JSON outputs.}
\label{tab:main-results}
\begin{tabular}{lccc}
\toprule
Configuration & Accuracy & Avg. thinking tokens & Notes \\
\midrule
Baseline ($n\_\text{wait}=0$) & TBA & TBA & pending execution \\
Budget forcing ($n\_\text{wait}=1$) & TBA & TBA & pending execution \\
Budget forcing ($n\_\text{wait}=2$) & TBA & TBA & pending execution \\
Budget forcing ($n\_\text{wait}=4$) & TBA & TBA & pending execution \\
\bottomrule
\end{tabular}
\end{table}

\subsection{Metrics}

The current evaluation code already computes scaling and performance directly from the accuracy sweep. Control may require additional explicit bookkeeping if target budgets are not recorded in the current output schema. Table~\ref{tab:metrics} should be populated from the final JSON artifact.

\begin{table}[H]
\centering
\caption{Control, scaling, and performance metrics for the final sweep.}
\label{tab:metrics}
\begin{tabular}{lc}
\toprule
Metric & Value \\
\midrule
Control & TBA \\
Scaling & TBA \\
Performance & TBA \\
\bottomrule
\end{tabular}
\end{table}

\subsection{Trigger Ablation}

The repository plan proposes an ablation over trigger phrases such as ``Wait'' and ``Hmm, let me reconsider''. This is a useful experiment because budget forcing assumes that the trigger phrase can coax the model into productive reflection instead of degenerate repetition. Table~\ref{tab:trigger} is reserved for that comparison.

\begin{table}[H]
\centering
\caption{Trigger phrase ablation.}
\label{tab:trigger}
\begin{tabular}{lcc}
\toprule
Trigger & Accuracy & Avg. thinking tokens \\
\midrule
Wait & TBA & TBA \\
Hmm, let me reconsider & TBA & TBA \\
\bottomrule
\end{tabular}
\end{table}

\subsection{Environment Log and Blockers}

This subsection should record the actual machine, package environment, model choice, and execution blockers encountered by Agent A. If the default configuration cannot run locally, that fact should be documented explicitly rather than hidden. A transparent blocker log improves the scientific value of the reproduction because it clarifies which claims were tested directly and which remained conditional on access to stronger hardware.

\subsection{Interim Interpretation}

Before final results are integrated, the most defensible interim conclusion is that the repository implements the intended experiment surface but has not yet completed the full evidence loop from code to saved artifact. The sprint goal is to close that gap by either producing the JSON outputs or documenting the exact operational reason why that could not be completed in the current environment.
\section{Discussion and Conclusion}\label{sec:conclusion}

This project is intentionally narrow. Rather than attempting to reproduce the entire s1 training recipe, it concentrates on the part of the method that is easiest to inspect and validate locally: budget forcing at decoding time. That focus is appropriate for a course project because it yields a useful test of reproducibility without turning the project into an infrastructure exercise.

The practical lesson is that simple inference-time interventions are worth studying precisely because they are falsifiable. If budget forcing improves accuracy in the local setup, the report can analyze how much improvement is obtained per unit of extra compute and whether the gain appears stable across triggers. If it fails or produces mixed results, that outcome is still valuable because it clarifies the dependency of the original claim on model family, prompt format, hardware, and benchmark details.

Future work should study three directions. First, trigger selection should move beyond fixed phrases toward adaptive or learned continuation prompts. Second, anti-loop safeguards should be incorporated so that extra thinking budget is spent on verification rather than repetition. Third, budget forcing should be tested in settings where reasoning must be combined with external knowledge or tools, including retrieval-augmented pipelines and non-mathematical tasks.
\label{sec:conclusion}

This repository captures an instructive slice of the broader reasoning-model landscape. Its main value is not that it recreates every part of the original s1 pipeline, but that it operationalizes one of the paper's clearest ideas: test-time compute can be controlled directly at decoding time. That makes budget forcing a useful teaching and reproduction target because the core mechanism is understandable, inspectable, and experimentally modifiable.

At the same time, the project also highlights a recurring tension in contemporary LLM research. Methods advertised as simple may still depend on heavyweight runtime conditions that are invisible in the algorithm sketch. For that reason, a careful reproduction should treat environment notes, model substitutions, and execution blockers as first-class results rather than incidental details.

The final conclusion of this report should be updated once Agent A finishes the runtime validation and result generation. If the experiments succeed, the discussion should explain whether additional enforced thinking improved accuracy and how stable that improvement was across triggers. If the experiments are blocked, the conclusion should instead emphasize what the repository confirms at the implementation level and what remains unverified at the systems level.

\bibliographystyle{apalike}
\bibliography{report/references}

\end{document}